
\medskip

%Dans ce  questionnaire à  choix multiples, pour chaque  question des réponses sont proposées, une seule est exacte. Sur la copie, écrire le numéro de la question et recopier la bonne réponse.
%
%\textbf{Pour la question 4, une justification est attendue}.
%
%\begin{center}
%\begin{tabularx}{\linewidth}{|m{6cm}|*{3}{>{\centering \arraybackslash}X|}}\hline
%\textbf{Questions}&A&B&C\\ \hline
%\textbf{1.} La décomposition en produit de facteurs premiers de 24 est : &$2  \times 3 \times 4$&$2\times 2\times 2\times 3$&$2\times 2\times 6$\\ \hline
%\textbf{2.} Lequel de ces  nombres est premier ? &\np{2255}&\np{8191}&\np{7113}\\ \hline
%\textbf{3.} La roue B fait 2 tours, combien
%de tours fait la roue A ?
%
%\includegraphics[width=3cm]{engrenage_Polynesie_2019} &3 tours&4 tours&5 tours\\ \hline
%\textbf{4.}  Pour cette question, une
%\textbf{justification est attendue}.
%
%\psset{unit=0.9cm}
%\begin{pspicture}(5.2,4.6)
%%\psgrid
%\psline(0,3.1)(4,4.3)
%\psline(0.6,0.1)(4.6,1.3)
%\psline(0,3.7)(3.9,0.5)
%\psline(2.1,0)(3.3,4.5)
%\uput[dl](2.2,0.6){P}\uput[u](2.5,1.63){R}\uput[u](0.5,3.25){T}\uput[ur](3.2,4){S}
%\uput[d](3.38,0.95){V}
%\rput{-40}(1.2,2.3){7,2~cm}\rput{14}(1.9,3.9){8,4~cm}
%\rput{-38}(3.2,1.4){3~cm}\rput(3.8,2.4){(TS) // (PV)}
%\end{pspicture}&PV = 3~cm&PV = 20,16~cm& PV = 3,5~cm\\
%\hline
%\end{tabularx}
%\end{center}
\begin{enumerate}
\item $24 = 8 \times 3 = 2^3 \times 3 = 2\times 2\times 2\times 3$ : réponse B.
\item \np{2255} est un multiple de 5 : il n'est pas premier.

La somme $7 + 1 + 1 + 3 = 12$ est un multiple de 3, donc \np{7113} est un multiple de 3 : il n'est pas premier. Il reste \np{8191} premier. Réponse B.
\item Les deux roues seront à nouveau en contact au même point qu'au départ quand les deux roues auront fait un nombre entiers de tours.

Les multiples de 12 sont : 12 ; 24 ; 36 ; 48 ; \ldots

Les multiples de 18 sont : 18 ; 36 ; 54 ; \ldots

On a donc $2 \times 18 = 3 \times 12$. Quand la roue B fait 2 tours, la roue A en fait 3. Réponse A.
\item Les droites (TS) et (PV) étant parallèles, on a une configuration de Thalès. On a donc :

$\dfrac{\text{RV}}{\text{RT}} = \dfrac{\text{PV}}{\text{ST}}$ soit $\dfrac{3}{7,2} = \dfrac{\text{PV}}{8,4}$, d'où PV $ = \dfrac{3}{7,2} \times 8,4 = \dfrac{25,2}{7,2} = 3,5$~(cm). Réponse C.
\end{enumerate}

\bigskip

